\documentclass[12pt,a4paper]{article}

\usepackage[a4paper, left=2.5cm, right=2.5cm, top=3cm, bottom=3cm]{geometry}  
\usepackage[brazil]{babel}  
\usepackage{float}  
\usepackage{setspace}       
\usepackage{graphicx}       
\usepackage{amsmath}        
\usepackage{hyperref}       
\usepackage{cleveref} 
\usepackage{lipsum}         
\usepackage{textcomp}       
\usepackage{booktabs}
\usepackage{csquotes}
\usepackage{array}     
\usepackage{amssymb}
\usepackage{bookmark} 
\usepackage[utf8]{inputenc}
\usepackage[date=year,
     backend=biber,
     backref=page,
     style=numeric-comp,
     bibstyle=ieee, 
     url=false,
     isbn=true,10.1109/TPWRS.2009.2037006
     sorting=nty,
     sortcites=true]{biblatex}
\usepackage{amsthm}
\newtheorem{definition}{Definição}[section]
\newtheorem{proposition}{Proposição}[section]

\usepackage{url}
\usepackage{xurl}
	 
\doublespacing

 \AtEveryBibitem{\clearfield{pages}}
 \AtEveryBibitem{\clearfield{notes}}

\addbibresource{ic.bib}

\title{\textbf{Árvores de Decisão}}
\author{Sofia Valverde Villas Bôas}
\date{}

\begin{document}

\maketitle

\section{Dúvidas minhas}

\begin{enumerate}
    \item Qual a melhor forma de eu traduzir splits?
\end{enumerate}

\section{Non‑greedy trees}

Árvores geradas da forma tradicional gulosa podem cair em ótimos locais, mas não necessariamente em globais. Isso porque o processo de construção dessas árvores se baseia em achar a melhor forma de dividir a partir de certo nó e não leva em consideração possíveis divisões futuras.

Existem algumas formas na literatura que sugerem como tratar esse problema de sub-otimalidade.

\subsection{Árvores Evolucionais}

Baseiam-se em \emph{algoritmos genéticos}




\begingroup
\onehalfspacing
\printbibliography
\endgroup

\end{document}
